\documentclass[10pt,a4paper]{amsart}
\usepackage[margin=19mm]{geometry}
\usepackage{amsmath,amssymb,mathtools}
\usepackage[scr=rsfs,cal=euler]{mathalfa}
\usepackage{xfrac}
\usepackage[unicode,hidelinks,bookmarks=false]{hyperref}
\newcommand{\ie}{i.e.\ }
\newcommand{\cf}{cf.\ }
\newcommand{\resp}{respectively\ }
\newcommand{\Subsection}{\subsection}
\providecommand{\nobreakdash}{\nobreak\hbox{-}}
\setlength{\emergencystretch}{2em}
\multlinegap0pt
\usepackage{bm}
\usepackage{bbm}
\usepackage{dsfont}
\usepackage{xspace}
\usepackage{cancel}
\usepackage{mathtools}
\usepackage[shortlabels]{enumitem}
\usepackage{amsthm}
\makeatletter
\@ifundefined{lemma}{\newtheorem{lemma}{Lemma}}{}
\@ifundefined{prop}{\newtheorem{prop}{Proposition}}{}
\makeatother
\title[On preconditioned Riemannian gradient methods for minimizing]{On preconditioned Riemannian gradient methods for minimizing the Gross-Pitaevskii energy functional: algorithms, global convergence and optimal local convergence rate}
\author{Zixu Feng, Qinglin Tang}
\date{Source: arXiv:2510.13516v3 (CC BY 4.0)}
\begin{document}
\maketitle

\noindent\emph{Source and licence.} On preconditioned Riemannian gradient methods for minimizing the Gross-Pitaevskii energy functional: algorithms, global convergence and optimal local convergence rate. Version arXiv:2510.13516v3 (CC BY 4.0), distributed under the Creative Commons Attribution 4.0 International licence, as recorded in the arXiv licence metadata for this exact version and shown on the abstract page https://arxiv.org/abs/2510.13516v3. Full rights notice: licenses/arxiv-2510.13516-CC-BY-4.0.txt.\par\smallskip

\noindent\textit{Editorial scope.} Counted unit: the Lemma whose source label is \texttt{Polyak-Loj} in the article (source0.tex lines 788--793, source label \texttt{Polyak-Loj}), delivered together with the article's own complete original proof of it (source0.tex lines 794--858, one \texttt{proof} environment of 65 lines). No mathematics is written, repaired or re-derived: statement and proof are the article's own text line for line. Required context retained uncounted: Property-of-E'' (lines 482--502); Property-P (lines 577--603); Transfer (lines 727--736). Every definition, display and label of those lines is retained unchanged. Omitted from this package and not redistributed: 1--481, 503--576, 604--726, 737--787, 859--2053. Nothing inside a retained excerpt is omitted or altered: every retained body is delivered exactly as the article's lines are written, so no omission receipt of any kind accompanies this package. Citations: the retained excerpts carry no citation command at all, so no bibliography entry of the article is transcribed and no reference is redistributed. Reference adaptation: none was needed and none was made; every reference inside the delivered unit resolves inside it, so no unresolved reference or citation is produced. Printed slips kept literal: no printed word, symbol, hypothesis or number of the counted statement or of its proof was altered. Layout and class: the article's own class is \texttt{elsarticle} with packages including stmaryrd, amsmath, cases, amssymb, amsfonts, bm, amsthm, dsfont, color, graphicx, caption, subcaption, xcolor, hyperref; the wrapper is the lane's pinned \texttt{amsart} editorial wrapper at 10pt on A4, which restates the theorem-like environments the retained text needs and adds the article's own macro definitions that the retained text needs (none), with extra wrapper packages (bm, bbm, dsfont, xspace, cancel, mathtools, enumitem). The delivered numbering is the unit's own. Assets: no retained span contains a figure environment, a graphics-inclusion command, a PDF-page inclusion, a table or a TikZ picture; no image file is copied into this package, the package declares no asset dependency and no image file is redistributed with it. Licence: version arXiv:2510.13516v3 of this article is distributed under the Creative Commons Attribution 4.0 International licence, as recorded in the arXiv licence metadata for this exact version and shown on the abstract page https://arxiv.org/abs/2510.13516v3. A full rights notice is delivered at licenses/arxiv-2510.13516-CC-BY-4.0.txt. Characters: every retained excerpt is pure ASCII, so the delivered engine, XeLaTeX with its default Latin Modern text font, reports no missing character.
\begin{prop}\label{Property-of-E''}
	Given $\phi\in H_0^1(\mathcal{D})$ and for all $u, v\in H_0^1(\mathcal{D})$, the following conclusions hold:
	\begin{itemize}[label={}, labelsep=0pt, leftmargin=*]
		\item $(i)$ {The second variation} $E''(\phi)$ satisfies the invariance under the following linear group actions
		\begin{align*}
			\left\langle E''(I_{\alpha}^{\beta}\phi)I_{\alpha}^{\beta}v,I_{\alpha}^{\beta}v\right\rangle=\left\langle E''(\phi)v,v\right\rangle \quad for\ all\ \alpha,\beta\in[-\pi, \pi). 
		\end{align*}
		\item $(ii)$ {The second variation} $E''(\phi)$ is a continuous operator on $H_0^1(\mathcal{D})$, i.e.,
		\begin{align*}
			\left|\big\langle E''(\phi)u,v\big\rangle\right|\le C_{\phi}\|u\|_{H^1}\|v\|_{H^1}.
		\end{align*}
		\item $(iii)$ Given $\psi\in H_0^1(\mathcal{D})$, for {$p_0=6/(4-\theta)\in\left[\frac{3}{2},6\right)$}, the following inequality holds
		\begin{align*}
			\left|\left\langle \big(E''(\phi)-E''(\psi)\big)u,v\right\rangle\right|\le C_{\phi,\psi}\|u\|_{H^1}\|v\|_{H^1}\|\phi-\psi\|_{L^{p_0}}.
		\end{align*}
		\item $(iv)$ The following Lipschitz-type inequality holds
		\begin{align*}
			E(\phi+v)-E(\phi)\le \big\langle E'(\phi),v\big\rangle+\frac{1}{2}\big\langle E''(\phi)v,v\big\rangle+C_{\phi,v}\|v\|^3_{H^1}.
		\end{align*}
	\end{itemize}
\end{prop}

\begin{prop}\label{Property-P}
	Assume {\rm \textbf{(A1)}}-{\rm \textbf{(A6)}}. Given $\phi\in H_0^1(\mathcal{D})$ and for all $u, v\in H_0^1(\mathcal{D})$ and $w\in H^{-1}(\mathcal{D})$, the following conclusions hold:
	\begin{itemize}[label={}, labelsep=0pt, leftmargin=*]
		\item $(i)$ If $E$ is a
		Morse-Bott functional on $\mathcal{S}$, then for all $\phi\in\mathcal{S}$, $\mathcal{P}_{\phi}$ and $E''(\phi)-\lambda_g\mathcal{I}$ satisfy the spectral equivalence, i.e.,
		\begin{align}\label{Mu-L}
			\inf_{v\in N_{\phi}\mathcal{M}\backslash\{0\}}\frac{\big\langle \big(E''(\phi)-\lambda_g\mathcal{I}\big)v,v\big\rangle}{\big\langle\mathcal{P}_{\phi}v,v\big\rangle}=\mu>0,\;\sup_{v\in T_{\phi}\mathcal{M}\backslash\{0\}}\frac{\big\langle \big(E''(\phi)-\lambda_g\mathcal{I}\big)v,v\big\rangle}{\big\langle\mathcal{P}_{\phi}v,v\big\rangle}=L<\infty.
		\end{align}
		\item $(ii)$ {The operator} $\mathcal{P}^{-1}_{\phi}\mathcal{H}_{\phi}:H_0^1(\mathcal{D})\to H_0^1(\mathcal{D})$ is a bounded linear operator, i.e., 
		\begin{align*}
			\big\|\mathcal{P}^{-1}_{\phi}\mathcal{H}_{\phi}v\big\|_{H^1}\le C_{\phi}\|v\|_{H^{1}}.
		\end{align*}
		Furthermore, $\mathcal{P}_{\phi}^{-1}(\mathcal{H}_{\phi}-\mathcal{P}_{\phi})$ satisfies the following estimate:
		\begin{align*}
			\big\|\mathcal{P}_{\phi}^{-1}(\mathcal{H}_{\phi}-\mathcal{P}_{\phi})v\big\|_{H^1}\le C_{\phi}\|v\|_{L^{p}}\quad\text{with}\quad {p=\max\{p_0,p_2\}\in[1,6)}.
		\end{align*}
		\item $(iii)$ Let $\phi\in\mathcal{M}$, there exists $\sigma$ such that for all $\psi\in\mathcal{B}_{\sigma}(\phi)$, the operator $\nabla^{\mathcal{R}}_{\mathcal{P}} E(\cdot):H_0^1(\mathcal{D})\to H_0^1(\mathcal{D})$ and the functional $\lambda_{(\cdot)}:H^1_0(\mathcal{D})\to \mathbb{R}$ are local Lipschitz continuous at $\phi$, i.e., 
		\begin{align*}
			\big\|\nabla^{\mathcal{R}}_{\mathcal{P}} E(\phi)-\nabla^{\mathcal{R}}_{\mathcal{P}} E(\psi)\big\|_{H^1}\le C_{\phi}\|\phi-\psi\|_{H^1}\quad and\quad \big|\lambda_{\phi}-\lambda_{\psi}\big|\le C_{\phi}\|\phi-\psi\|_{L^p},
		\end{align*}
		where $p=\max\{p_0,p_1,p_2,2\}\in[1,6)$.
		\item $(iv)$ Let $\phi\in\mathcal{M}$, for all $v\in T_{\phi}\mathcal{M}$, it holds that
		\begin{align*}
			\big|\mathfrak{R}_{\phi}(tv)-(\phi+tv)\big|\le\frac{1}{2}t^2\|v\|^2_{L^2}|\phi+tv|.
		\end{align*}
	\end{itemize}
\end{prop}

\begin{lemma}\label{Transfer}
	Let $E$ be a
	Morse-Bott functional on $\mathcal{S}$. For any $\phi \in \mathcal{M}$ and $\phi_g\in\mathcal{S}$, there exists $\phi^*_g \in \mathcal{S}$ such that the following orthogonality conditions hold:
	{\rm
		\begin{align*}
			(\phi-\phi_g^*,i\phi_g^*)_{L^2}=0\quad and\quad(\phi-\phi_g^*,i\mathcal{L}_z\phi_g^*)_{L^2}=0.
		\end{align*}
	}
	Furthermore, $\|\phi-\phi_g^*\|_{H^1}\le C_{\phi} \|\phi-\phi_g\|_{H^1}.$
\end{lemma}

\begin{lemma}\label{Polyak-Loj}
	Let $E$ be a Morse-Bott functional on $\mathcal{S}$. For any $\phi_g\in\mathcal{S}$, and for every sufficiently small $\varepsilon>0$, there exists $\sigma>0$ such that for any $\phi \in \mathcal{B}_{\sigma}(\phi_g)$, the following Polyak-\L ojasiewicz inequality holds:
	\[
	E(\phi) - E(\phi_g)\leq \frac{1}{2(\mu - \varepsilon)} \left\| \nabla^{\mathcal{R}}_{\mathcal{P}} E(\phi)\right\|^2_{\mathcal{P}_{\phi}}.
	\]
\end{lemma}

\begin{proof}
	According to $E(\phi_g^*)=E(\phi_g)$ and Taylor's formula at $\phi$, we have
	\begin{align}
		\nonumber E(\phi)-&E(\phi_g)=E(\phi)-E(\phi_g^*)\\
		\nonumber=&\left\langle E'(\phi),\phi-\phi_g^*\right\rangle-\frac{1}{2}\left\langle E''(\phi)(\phi-\phi_g^*),\phi-\phi_g^*\right\rangle+o\left(\|\phi-\phi_g^*\|^2_{H^1}\right)\\
		=&\left( \nabla^{\mathcal{R}}_{\mathcal{P}} E(\phi),\phi-\phi_g^*\right)_{\mathcal{P}_{\phi}}\hspace{-0.32cm}-\frac{1}{2}\left\langle \big(E''(\phi)-\lambda_{\phi}\mathcal{I}\big)(\phi-\phi_g^*),\phi-\phi_g^*\right\rangle+o\left(\|\phi-\phi_g^*\|^2_{H^1}\right)\label{P-L0}.
	\end{align}
	Note that
	\begin{align}
		\nonumber\phi-\phi_g^*&=-\phi_g^*+(\phi_g^*,\phi)_{L^2}\phi+(\phi-\phi_g^*,\phi)_{L^2}\phi\\
		\nonumber&=\phi-\phi_g^*+(\phi_g^*-\phi,\phi)_{L^2}\phi+\frac{1}{2}\left(\|\phi\|_{L^2}^2-\|\phi_g^*\|_{L^2}^2+\|\phi-\phi_g^*\|_{L^2}\right)\phi\\
		&=\text{Proj}^{L^2}_{\phi}(\phi-\phi_g^*)+\frac{1}{2}\|\phi-\phi_g^*\|^2_{L^2}\phi,\label{P-L1-1}\\
		\nonumber\phi-\phi_g^*&=\phi-(\phi,\phi_g^*)_{L^2}\phi_g^*-(\phi_g^*-\phi,\phi_g^*)_{L^2}\phi_g^*\\
		&=\text{Proj}^{L^2}_{\phi_g^*}(\phi-\phi_g^*)-\frac{1}{2}\|\phi-\phi_g^*\|^2_{L^2}\phi_g^*,\label{P-L1}\\
		\Longrightarrow\;&\text{Proj}^{L^2}_{\phi}(\phi-\phi_g^*)=\text{Proj}^{L^2}_{\phi_g^*}(\phi-\phi_g^*)-\frac{1}{2}\|\phi-\phi_g^*\|^2_{L^2}\phi-\frac{1}{2}\|\phi-\phi_g^*\|^2_{L^2}\phi_g^*,\label{P-L1-2}
	\end{align}
	where $\text{Proj}^{L^2}_{\phi_g^*}(\phi-\phi_g^*)\in N_{\phi_g^*}\mathcal{M}$.
	Substituting \eqref{P-L1-1} into \eqref{P-L0}, and using {\bf Proposition \ref{Property-of-E''}}-$(ii)$ and {\bf Proposition \ref{Property-P}}-$(iii)$, we derive
	\begin{align*}
		E(\phi)-E(\phi_g)&=\left( \nabla^{\mathcal{R}}_{\mathcal{P}} E(\phi),\text{Proj}^{L^2}_{\phi}(\phi-\phi_g^*)\right)_{\mathcal{P}_{\phi}}\\
		\nonumber-&\frac{1}{2}\left\langle \big(E''(\phi)-\lambda_{\phi}\mathcal{I}\big)\text{Proj}^{L^2}_{\phi}(\phi-\phi_g^*),\text{Proj}^{L^2}_{\phi}(\phi-\phi_g^*)\right\rangle+o\left(\|\phi-\phi_g^*\|^2_{H^1}\right).
	\end{align*}
	Plugging \eqref{P-L1-2} into the above identity, we get
	
	\begin{align}\label{P-L2}
		\nonumber E(\phi)-E(\phi_g)&=\left(\nabla^{\mathcal{R}}_{\mathcal{P}} E(\phi),\text{Proj}^{L^2}_{\phi_g^*}(\phi-\phi_g^*)\right)_{\mathcal{P}_{\phi}}\\
		-&\frac{1}{2}\left\langle \big(E''(\phi)-\lambda_{\phi}\mathcal{I}\big)\text{Proj}^{L^2}_{\phi_g^*}(\phi-\phi_g^*),\text{Proj}^{L^2}_{\phi_g^*}(\phi-\phi_g^*)\right\rangle+o\left(\|\phi-\phi_g^*\|^2_{H^1}\right).
	\end{align}
	Based on {\bf Proposition \ref{Property-of-E''}}-$(iii)$, {\bf Proposition \ref{Property-P}}-$(iii)$, and {\bf (A6)}-$(iii)$, the following estimations hold
	\begin{align*}
		\left\langle \big(E''(\phi)-E''(\phi_g^*)\big)\text{Proj}^{L^2}_{\phi_g^*}(\phi-\phi_g^*),\text{Proj}^{L^2}_{\phi_g^*}(\phi-\phi_g^*)\right\rangle&=o\left(\|\phi-\phi_g^*\|^2_{H^1}\right),\\
		\left\langle \big(\lambda_{\phi_g^*}\mathcal{I}-\lambda_{\phi}\mathcal{I}\big)\text{Proj}^{L^2}_{\phi_g^*}(\phi-\phi_g^*),\text{Proj}^{L^2}_{\phi_g^*}(\phi-\phi_g^*)\right\rangle&=o\left(\|\phi-\phi_g^*\|^2_{H^1}\right),\\
		\left\langle\big(\mathcal{P}_{\phi}-\mathcal{P}_{\phi_g^*}\big)\text{Proj}^{L^2}_{\phi_g^*}(\phi-\phi_g^*),\text{Proj}^{L^2}_{\phi_g^*}(\phi-\phi_g^*)\right\rangle&=o\left(\|\phi-\phi_g^*\|^2_{H^1}\right).
	\end{align*}
	According to {\bf Proposition \ref{Property-P}}-$(i)$ and $\text{Proj}^{L^2}_{\phi_g^*}(\phi-\phi_g^*)\in N_{\phi_g^*}\mathcal{M}$, the following lower bound estimate holds
	\begin{align*}
		\frac{\left\langle\big( E''(\phi_g^*)-\lambda_{g}\mathcal{I}\big)\text{Proj}^{L^2}_{\phi_g^*}(\phi-\phi_g^*),\text{Proj}^{L^2}_{\phi_g^*}(\phi-\phi_g^*)\right\rangle}{\left\langle\mathcal{P}_{\phi_g^*}\text{Proj}^{L^2}_{\phi_g^*}(\phi-\phi_g^*),\text{Proj}^{L^2}_{\phi_g^*}(\phi-\phi_g^*)\right\rangle}\ge\mu.
	\end{align*}
	In summary, the estimate we want is derived
	\begin{align*}
		-\frac{1}{2}&\left\langle \big(E''(\phi)-\lambda_{\phi}\mathcal{I}\big)\text{Proj}^{L^2}_{\phi_g^*}(\phi-\phi_g^*),\text{Proj}^{L^2}_{\phi_g^*}(\phi-\phi_g^*)\right\rangle\\
		&\qquad\qquad\qquad\le-\frac{\mu}{2}\left\langle\mathcal{P}_{\phi}\text{Proj}^{L^2}_{\phi_g^*}(\phi-\phi_g^*),\text{Proj}^{L^2}_{\phi_g^*}(\phi-\phi_g^*)\right\rangle
		+o\left(\|\phi-\phi_g^*\|^2_{H^1}\right).
	\end{align*}
	Combining the above inequality with \eqref{P-L2}, we get
	\begin{align}\label{P-L3}
		\nonumber E(\phi)-E(\phi_g)&\le\left(\nabla^{\mathcal{R}}_{\mathcal{P}} E(\phi),\text{Proj}^{L^2}_{\phi_g^*}(\phi-\phi_g^*)\right)_{\mathcal{P}_{\phi}}\\
		&-\frac{\mu}{2}\left(\text{Proj}^{L^2}_{\phi_g^*}(\phi-\phi_g^*),\text{Proj}^{L^2}_{\phi_g^*}(\phi-\phi_g^*)\right)_{\mathcal{P}_{\phi}}+o\left(\|\phi-\phi_g^*\|^2_{H^1}\right).
	\end{align}
	By {\bf Lemma \ref{Transfer}} and {\bf (A6)}-$(ii)$, we know that 
	\begin{align}\label{phi_g^*-phi_g}
		\|\phi-\phi_g^*\|_{{H^1}}\le C\|\phi-\phi_g^*\|_{\mathcal{P}_{\phi}}\le C_{\phi}\|\phi-\phi_g^*\|_{H^1}\le C_{\phi} \|\phi-\phi_g\|_{H^1}. 
	\end{align}
	Recalling \eqref{P-L1}, then for all sufficiently small $\varepsilon$, there exists $\sigma$ such that for any $\phi \in \mathcal{B}_{\sigma}(\phi_g)$, we have
	\begin{align}\label{P-L4}
		\left|o\left(\|\phi-\phi_g^*\|^2_{H^1}\right)\right|\le\frac{\varepsilon}{2}\left\|\text{Proj}^{L^2}_{\phi_g^*}(\phi-\phi_g^*)\right\|^2_{\mathcal{P}_{\phi}}.
	\end{align}
	Then, by \eqref{P-L3}, the Polyak-\L ojasiewicz inequality is deduced as follows
	\begin{align*}
		&E(\phi)-E(\phi_g)\\
		\le&\left(\nabla^{\mathcal{R}}_{\mathcal{P}} E(\phi),\text{Proj}^{L^2}_{\phi_g^*}(\phi-\phi_g^*)\right)_{\mathcal{P}_{\phi}}-\frac{\mu-\varepsilon}{2}\left(\text{Proj}^{L^2}_{\phi_g^*}(\phi-\phi_g^*),\text{Proj}^{L^2}_{\phi_g^*}(\phi-\phi_g^*)\right)_{\mathcal{P}_{\phi}}\\
		\le&\sup\limits_{v\in H_0^1(\mathcal{D})}\Bigg(\left( \nabla^{\mathcal{R}}_{\mathcal{P}} E(\phi),v\right)_{\mathcal{P}_{\phi}}-\frac{\mu-\varepsilon}{2}(v,v)_{\mathcal{P}_{\phi}}\Bigg)
		=\frac{1}{2(\mu-\varepsilon)}\left\|\nabla^{\mathcal{R}}_{\mathcal{P}} E(\phi)\right\|^2_{\mathcal{P}_{\phi}}.
	\end{align*}
\end{proof}

\end{document}
